By Theorem \claimref{singularity.rate-discharge-hypothesis-iii}, under (P, \claimref{singularity.p-univalence}) and (R, \claimref{singularity.r-regularity-fixed-locus}) the marginal of a mode of weight $\omega$ is $\operatorname{diag}(1 - n, n)$ with $n = 1/(e^{p\omega} + 1)$, $p = 2\pi/\kappa$; write $x = p\omega$. Its entropy is $s = -n\ln n - (1 - n)\ln(1 - n)$. With $r = e^{x}$, $n = 1/(r + 1)$ and $1 - n = r/(r + 1)$, so $s = \ln(r + 1) - \frac{r}{r + 1}\ln r = \ln(1 + e^{-x}) + x\,n$, which is the stated form. Exchanging $n$ and $1 - n$ sends $x$ to $-x$ and leaves $s$ unchanged, so $s$ is even. Differentiating, $\frac{dn}{dx} = -\frac{r}{(r + 1)^{2}} = -n(1 - n)$, and
\[
s'(x) \;=\; \frac{-e^{-x}}{1 + e^{-x}} + n + x\,\frac{dn}{dx} \;=\; -n + n - x\,n(1 - n) \;=\; -x\,n(1 - n) ,
\]
negative for $x > 0$ and positive for $x < 0$, since $0 < n < 1$. So $s$ has its unique maximum at $x = 0$, where $n = \tfrac12$ and $s = \ln 2$, and $s \to 0$ as $|x| \to \infty$ because $n \to 0$ or $1$. At $x = 0$ the gluing of (R) reads $n = 1 - n$ for every period, as the theorem records for zero-weight modes. The value at $x = 2\pi$ follows from $e^{-2\pi} = 0.0018674\ldots$: $s(2\pi) = \ln(1.0018674\ldots) + 2\pi \times 0.0018639\ldots = 0.01358\ldots$, against $\ln 2 = 0.6931\ldots$.
